\SourcePage{446}
\SGAsection{3}{爆破概形的 $K'$ 的计算}
\textbf{3.1.} 本节以 $i:Y\to X$ 表示一个\Term{regular-closed-immersion}{正则闭浸入}，以 $X'$ 表示沿 $Y$ \Term{blow-up}{爆破}所得的\Term{scheme}{概形}，并置 $Y'=Y\times_X X'$。
\SourcePage{447}
由 EGA II 8.1.8，$Y'$ 是 $X'$ 的一个\Term{closed-subscheme}{闭子概形}，由 $\mathcal O_{X'}$ 的一个\Term{ideal}{理想}定义，此\Term{ideal}{理想}\Term{canonical-isomorphism}{典范同构}于 $\mathcal O_{X'}(1)$；$j:Y'\to X'$ 由一个\Term{invertible-ideal}{可逆理想}定义，因而是\Term{codimension}{余维数}为 $1$ 的\Term{regular-closed-immersion}{正则闭浸入}。以 $\mathcal J$ 表示定义 $Y$ 的 $\mathcal O_X$ 的\Term{ideal}{理想}，以 $\mathcal J'$ 表示定义 $Y'$ 的 $\mathcal O_{X'}$ 的\Term{ideal}{理想}；于是
\[
\mathcal J'\simeq\mathcal J\mathcal O_{X'}\simeq\mathcal O_{X'}(1).
\]
最后，以 $f:X'\to X$ 和 $g:Y'\to Y$ 表示\Term{structure-morphism}{结构态射}。我们有\Term{commutative-diagram}{交换图}
\[
\begin{tikzcd}
Y'\arrow[r,"j"]\arrow[d,"g"']&X'\arrow[d,"f"]\\
Y\arrow[r,"i"']&X
\end{tikzcd}\tag{3.1.1}
\]
由定义，$X'=\operatorname{Proj}(\mathcal S)$，其中 $\mathcal S=\bigoplus_{n\geq0}\mathcal J^n$。因此
\[
Y'=\operatorname{Proj}(i^*\mathcal S)
=\operatorname{Proj}\left(\left(\bigoplus_{n\geq0}\mathcal J^n\right)\otimes_{\mathcal O_X}(\mathcal O_X/\mathcal J)\right)
=\operatorname{Proj}\left(\bigoplus_{n\geq0}\mathcal J^n/\mathcal J^{n+1}\right).
\]
由于 $i$ \Term{regular-immersion}{正则}，由 1.3，$\bigoplus_{n\geq0}\mathcal J^n/\mathcal J^{n+1}\simeq\operatorname{Sym}_{\mathcal O_Y}(\mathcal J/\mathcal J^2)$。以 $\mathcal N$ 表示\Term{conormal-sheaf}{余法层} $\mathcal J/\mathcal J^2$，便有
\[
Y'=\mathbb P(\mathcal N).
\]
在 $Y'$ 上有\Term{exact-sequence}{正合列}
\[
0\longrightarrow\mathcal F\longrightarrow g^*(\mathcal N)\longrightarrow\mathcal O_{Y'}(1)\longrightarrow0.\tag{3.1.2}
\]
注意，$\mathcal O_{Y'}(1)=j^*(\mathcal O_{X'}(1))=\mathcal J'/\mathcal J'^2$，故它正是\Term{immersion}{浸入} $j$ 的\Term{conormal-sheaf}{余法层}；以下记作 $\mathcal L$。

\begin{SGAresult}{\Term{lemma}{引理}}{3.2}
在 3.1 的假设和记号下，有如下\Term{isomorphism}{同构}：
\[
\begin{aligned}
\text{i)}\quad&\text{对每个 }i\geq0,\quad
\operatorname{Tor}^{\mathcal O_X}_i(\mathcal O_{X'},\mathcal O_Y)
\simeq\bigwedge^i\operatorname{Tor}^{\mathcal O_X}_1(\mathcal O_{X'},\mathcal O_Y);\\
\text{ii)}\quad&\operatorname{Tor}^{\mathcal O_X}_1(\mathcal O_{X'},\mathcal O_Y)\simeq\mathcal F.
\end{aligned}
\]
\end{SGAresult}
\begin{proof}
\SourcePage{448}
i) 因为 $\mathcal O_Y$ \Term{local-scheme}{局部}可用\Term{koszul-complex}{Koszul 复形}\Term{resolution}{解消}，\Term{graded-algebra}{分次代数} $\bigoplus_{i\geq0}\operatorname{Tor}^{\mathcal O_X}_i(\mathcal O_{X'},\mathcal O_Y)$ 是交替的（2.2）。因此存在唯一的\Term{algebra-homomorphism}{代数同态}
\[
\bigwedge\operatorname{Tor}^{\mathcal O_X}_1(\mathcal O_{X'},\mathcal O_Y)
\longrightarrow\bigoplus_{i\geq0}\operatorname{Tor}^{\mathcal O_X}_i(\mathcal O_{X'},\mathcal O_Y),\tag{3.2.1}
\]
延拓 $\operatorname{Tor}^{\mathcal O_X}_1(\mathcal O_{X'},\mathcal O_Y)$ 到上述\Term{direct-sum}{直和}中的\Term{injection}{嵌入}。只须\Term{local-scheme}{局部}证明它是\Term{isomorphism}{同构}。

所以可在 $X$ 的一个\Term{affine-open}{仿射开集} $U=\operatorname{Spec}(A)$ 上工作，设 $\mathcal J$ 由 $A$ 的\Term{ideal}{理想} $J$ 定义，并给定一个定义 $J$ 的\Term{regular-homomorphism}{正则同态} $u:E\to A$，其中 $E=A^n$。设 $f_1,\ldots,f_d$ 是 $E$ 的一组\Term{basis}{基}的像。在 $\bigoplus_{n\geq0}J^n$ 中，对于 $m\leq p$，以 $x^{(m)}$ 表示将 $J^p$ 的元素 $x$ 看作\Term{degree}{次数}为 $m$ 的元素。可以进一步\Term{restriction}{限制}到 $X'$ 中 $f_1^{(1)}$ \Term{invertible}{可逆}的\Term{open-subset}{开集} $V$ 上。

各 $\operatorname{Tor}^{\mathcal O_X}_i(\mathcal O_{X'},\mathcal O_Y)$ 是\Term{complex}{复形} $K_\bullet(u)\otimes_{\mathcal O_X}\mathcal O_{X'}$ 的\Term{homology-sheaf}{同调层}。为在 $V$ 上研究此\Term{complex}{复形}，在 $E'=E\otimes_A\Gamma(V,\mathcal O_{X'})$ 中作如下\Term{base-change}{换基}：若 $e_1,\ldots,e_d$ 是 $E$ 原来的\Term{basis}{基}，置
\[
e'_1=e_1,\qquad e'_i=e_i-\frac{f_i^{(1)}}{f_1^{(1)}}e_1\quad(i\ne1).
\]
置 $A'=\Gamma(V,\mathcal O_{X'})$；\Term{complex}{复形} $K_\bullet(u)\otimes_A A'$ 的\Term{differential}{微分}为
\[
d(e'_1)=f_1^{(0)},\qquad d(e'_i)=0\quad(i\ne1).
\]
以 $e'_{i_1}\wedge\cdots\wedge e'_{i_k}$（$i_1<\cdots<i_k$）作为 $\bigwedge^k(E')$ 的\Term{basis}{基}，得到
\[
\begin{aligned}
d(ae'_{i_1}\wedge\cdots\wedge e'_{i_k})&=0&&\text{若 }i_1\ne1,\\
d(ae'_{i_1}\wedge\cdots\wedge e'_{i_k})&=af_1^{(0)}e'_{i_2}\wedge\cdots\wedge e'_{i_k}&&\text{若 }i_1=1.
\end{aligned}
\]
\SourcePage{449}
以 $E''$ 表示 $E'$ 中由 $e'_2,\ldots,e'_d$ 生成的\Term{submodule}{子模}，可见 $d:\bigwedge^k E'\to\bigwedge^{k-1}E'$ 的像为 $f_1^{(0)}\bigwedge^{k-1}E''$。另一方面，$f_1^{(0)}$ 在 $A'$ 中\Term{regular-element}{正则}，因为在 $V$ 上\Term{invertible-ideal}{可逆理想} $\mathcal J'$ 由方程 $f_1^{(0)}$ 定义。因此 $d:\bigwedge^kE'\to\bigwedge^{k-1}E'$ 的\Term{kernel}{核}为 $\bigwedge^kE''$。所以 $K_\bullet(u)\otimes_AA'$ 的第 $k$ 个\Term{homology-group}{同调群}是 $\bigwedge^kE''/f_1^{(0)}\bigwedge^kE''$，在 $V$ 上便得到\Term{isomorphism}{同构}
\[
\operatorname{Tor}^{\mathcal O_X}_i(\mathcal O_{X'},\mathcal O_Y)
\simeq(\mathcal O_{X'}/f_1^{(0)}\mathcal O_{X'})\otimes_{\mathcal O_{X'}}\bigwedge^i(\mathcal O_{X'})^{d-1}
\simeq\bigwedge^i(\mathcal O_{Y'})^{d-1}.\tag{3.2.2}
\]
特别地，$i=1$ 时得到
\[
\operatorname{Tor}^{\mathcal O_X}_1(\mathcal O_{X'},\mathcal O_Y)\simeq(\mathcal O_{Y'})^{d-1}.
\]
于是对每个 $i$ 都有
\[
\operatorname{Tor}^{\mathcal O_X}_i(\mathcal O_{X'},\mathcal O_Y)
\simeq\bigwedge^i\operatorname{Tor}^{\mathcal O_X}_1(\mathcal O_{X'},\mathcal O_Y).\tag{3.2.3}
\]
由 2.2，(3.2.2) 所定义的\Term{isomorphism}{同构}
\[
\bigwedge(\mathcal O_{Y'})^{d-1}\xrightarrow{\sim}\bigoplus_{i\geq0}\operatorname{Tor}^{\mathcal O_X}_i(\mathcal O_{X'},\mathcal O_Y)
\]
是\Term{algebra-isomorphism}{代数同构}。所以 (3.2.3) 所定义的\Term{isomorphism}{同构}
\[
\bigwedge\operatorname{Tor}^{\mathcal O_X}_1(\mathcal O_{X'},\mathcal O_Y)
\xrightarrow{\sim}\bigoplus_{i\geq0}\operatorname{Tor}^{\mathcal O_X}_i(\mathcal O_{X'},\mathcal O_Y)
\]
也是\Term{algebra-isomorphism}{代数同构}。由唯一性，它就是\Term{homomorphism}{同态} (3.2.1) 在 $V$ 上的\Term{restriction}{限制}，故后者是\Term{isomorphism}{同构}。

ii) 考虑\Term{exact-sequence}{正合列}
\[
0\longrightarrow\mathcal J\longrightarrow\mathcal O_X\longrightarrow\mathcal O_Y\longrightarrow0,
\]
\SourcePage{450}
以及与 $\mathcal O_{X'}$ 作\Term{tensor-product}{张量积}所得的 Tor \Term{exact-sequence}{正合列}：
\[
0\longrightarrow\operatorname{Tor}^{\mathcal O_X}_1(\mathcal O_{X'},\mathcal O_Y)
\longrightarrow\mathcal J\otimes_{\mathcal O_X}\mathcal O_{X'}\longrightarrow\mathcal O_{X'}\longrightarrow\mathcal O_{Y'}\longrightarrow0.
\]
\Term{morphism}{态射} $\mathcal J\otimes_{\mathcal O_X}\mathcal O_{X'}\to\mathcal O_{X'}$ 的像为 $\mathcal J\mathcal O_{X'}=\mathcal O_{X'}(1)$，故有\Term{exact-sequence}{正合列}
\[
0\longrightarrow\operatorname{Tor}^{\mathcal O_X}_1(\mathcal O_{X'},\mathcal O_Y)
\longrightarrow\mathcal J\otimes_{\mathcal O_X}\mathcal O_{X'}\longrightarrow\mathcal O_{X'}(1)\longrightarrow0.
\]
再与 $\mathcal O_{Y'}$ 作\Term{tensor-product}{张量积}；因 $\operatorname{Tor}^{\mathcal O_X}_1(\mathcal O_{X'},\mathcal O_Y)$ 是 $\mathcal O_{Y'}$-\Term{module}{模}，且 $\mathcal O_{X'}(1)$ \Term{invertible-sheaf}{可逆}，得到\Term{exact-sequence}{正合列}
\[
0\longrightarrow\operatorname{Tor}^{\mathcal O_X}_1(\mathcal O_{X'},\mathcal O_Y)
\longrightarrow g^*(\mathcal N)\longrightarrow\mathcal O_{Y'}(1)\longrightarrow0.
\]
容易看出，其中的\Term{morphism}{态射} $g^*(\mathcal N)\to\mathcal O_{Y'}(1)$ 是\Term{canonical-map}{典范态射}，所以此列就是 (3.1.2)。由定义即有
\[
\operatorname{Tor}^{\mathcal O_X}_1(\mathcal O_{X'},\mathcal O_Y)=\mathcal F.
\]
\end{proof}

\textbf{3.3.} 现在假定 $X$ 是有\Term{ample-sheaf}{丰沛层}的\Term{scheme}{概形}。则 $X'$ 在 $X$ 上\Term{projective-morphism}{射影}，故 $X'$ 也有\Term{ample-sheaf}{丰沛层}，$Y$ 和 $Y'$ 亦然。此外，$i$ 和 $j$ 是\Term{regular-immersion}{正则浸入}，所以是\Term{perfect-morphism}{完美态射}；$g$ 是\Term{smooth-morphism}{光滑态射}，所以也是\Term{perfect-morphism}{完美态射}，因为 $Y'=\mathbb P(\mathcal N)$；最后，由 1.9，$f$ 也是\Term{perfect-morphism}{完美态射}。这些\Term{morphism}{态射}都是\Term{projective-morphism}{射影}的，因此可以定义如下\Term{additive-homomorphism}{加法同态}（第 IV 讲\SGAsourcegap{原文此处未填第 IV 讲的具体引用编号}）：
\[
\begin{aligned}
i_* &:K'(Y)\longrightarrow K'(X),&j_* &:K'(Y')\longrightarrow K'(X'),\\
f_* &:K'(X')\longrightarrow K'(X),&g_* &:K'(Y')\longrightarrow K'(Y).
\end{aligned}
\]
\SourcePage{451}
\begin{SGAresult}{命题}{3.4}
在 3.1 和 3.3 的假设与记号下，对每个 $x\in K'(Y)$ 都有
\[
f^*(i_*(x))=j_*(g^*(x)\lambda_{-1}(F)).
\]
\end{SGAresult}
\begin{proof}
待证等式两边都是可加的，因此可设 $x$ 是\Term{finite-locally-free}{有限秩局部自由} $\mathcal O_Y$-\Term{module}{模} $\mathcal M$ 的类。此时 $f^*(i_*(x))$ 就是\Term{complex}{复形} $\mathbf Lf^*(\mathbf Ri_*(\mathcal M))$ 的类；此\Term{complex}{复形}由 $i_*(\mathcal M)$ 的一个\Term{projective-resolution}{投射解消}经\Term{functor}{函子} $f^*$ 得到。所得\Term{complex}{复形}的\Term{homology-sheaf}{同调层}为 $\operatorname{Tor}^{\mathcal O_X}_i(\mathcal M,\mathcal O_{X'})$。由 2.4 和 3.2，
\[
\operatorname{Tor}^{\mathcal O_X}_i(\mathcal M,\mathcal O_{X'})
\simeq\mathcal M\otimes_{\mathcal O_Y}\operatorname{Tor}^{\mathcal O_X}_i(\mathcal O_{X'},\mathcal O_Y)
\simeq\mathcal M\otimes_{\mathcal O_Y}\bigwedge^i\mathcal F.
\]
所以这些\Term{homology-sheaf}{同调层}在 $\mathcal O_{Y'}$ 上\Term{finite-locally-free}{有限秩局部自由}，因而在 $\mathcal O_{X'}$ 上\Term{perfect-complex}{完美}。于是 $\mathbf Lf^*(\mathbf Ri_*(\mathcal M))$ 在 $K'(X')$ 中的类是其\Term{homology-sheaf}{同调层}的类的\Term{alternating-sum}{交错和}，即
\[
f^*(i_*(x))=\sum_{i\geq0}(-1)^i\operatorname{cl}\left(j_*\left(\mathcal M\otimes_{\mathcal O_Y}\bigwedge^i\mathcal F\right)\right)
=\sum_{i\geq0}(-1)^i j_*(g^*(x)\lambda^i(F)).
\]
\end{proof}

\begin{SGAresult}{\Term{lemma}{引理}}{3.5}
在 3.1 的假设和记号下，有如下\Term{isomorphism}{同构}：

 i) 对每个 $n\geq0$ 和每个 $i\geq1$，$R^if_*(\mathcal O_{X'}(n))=0$；

 ii) 对每个 $n\geq0$，$f_*(\mathcal O_{X'}(n))=\mathcal J^n$。
\end{SGAresult}
\begin{proof}
显然可设 $X$ \Term{affine}{仿射}，因而\Term{quasi-compact}{拟紧}。

i) 先回顾以下\Term{lemma}{引理}，它已在第\SGAsourcegap{原文“已在第……讲中证明”未填讲号及引用位置}讲中证明。
\SourcePage{452}
\begin{SGAresult}{\Term{lemma}{引理}}{}
设 $f:X'\to X$ 为\Term{projective-morphism}{射影}\Term{perfect-morphism}{完美态射}，$X$ \Term{quasi-compact}{拟紧}；设 $\mathcal L^\bullet$ 是 $\mathcal O_{X'}$ 上的一个\Term{perfect-complex}{完美复形}，在\Term{degree}{次数} $>0$ 无\Term{cohomology}{上同调}。则存在整数 $n_0$，使得当 $n\geq n_0$ 时，\Term{complex}{复形} $\mathbf Rf_*(\mathcal L^\bullet\otimes_{\mathcal O_{X'}}\mathcal O_{X'}(n))$ 在\Term{degree}{次数} $>0$ 无\Term{cohomology}{上同调}。
\end{SGAresult}
将此\Term{lemma}{引理}用于只含 $\mathcal O_{X'}$ 的\Term{complex}{复形}，得到整数 $n_0$，使 $n\geq n_0$ 时，对每个 $i\geq1$ 都有 $R^if_*(\mathcal O_{X'}(n))=0$。再作\Term{descent}{下降}归纳，直至 $n=0$。为此考虑\Term{exact-sequence}{正合列}
\[
0\longrightarrow\mathcal O_{X'}(n+1)\longrightarrow\mathcal O_{X'}(n)\longrightarrow\mathcal O_{Y'}(n)\longrightarrow0,
\]
并施以\Term{functor}{函子} $f_*$。由于 $\mathbf Rf_*(\mathcal O_{Y'}(n))=i_*\mathbf Rg_*(\mathcal O_{Y'}(n))$，且 $Y'=\mathbb P(\mathcal N)$，所以 $n\geq0$ 时 $\mathcal O_{Y'}(n)$ 是 $g_*$-\Term{acyclic-object}{无上同调的}。因此对每个 $n\geq0$ 和 $i\geq1$ 都有
\[
R^if_*(\mathcal O_{Y'}(n))=0,
\]
从而得到\Term{exact-sequence}{正合列}
\[
R^if_*(\mathcal O_{X'}(n+1))\longrightarrow R^if_*(\mathcal O_{X'}(n))\longrightarrow0,
\]
这就完成了归纳。

ii) 要证明\Term{canonical-map}{典范态射}
\[
\bigoplus_{n\geq0}\mathcal J^n\longrightarrow\bigoplus_{n\geq0}f_*(\mathcal O_{X'}(n))
\]
是\Term{isomorphism}{同构}。此性质是\Term{local-scheme}{局部}的，因此可设 $\mathcal J$ 由 $\mathcal O_X$ 的一族\Term{section}{截面} $f_1,\ldots,f_d$ 定义，且相应的\Term{homomorphism}{同态} $(\mathcal O_X)^d\to\mathcal O_X$ \Term{regular-homomorphism}{正则}。置 $\mathcal S=\bigoplus_{n\geq0}\mathcal J^n$，定义\Term{surjective}{满}的\Term{algebra-homomorphism}{代数同态}
\[
v:\mathcal O_X[T_1,\ldots,T_d]\longrightarrow\mathcal S,
\]
\SourcePage{453}
将 $T_i$ 映到 $f_i^{(1)}$。这个\Term{homomorphism}{同态}定义一个\Term{closed-immersion}{闭浸入} $X'\to\mathbb P_X^{d-1}$，由 1.8 ii)，它是\Term{regular-immersion}{正则}的。

由 1.8 iii)，$\mathcal S$ 作为\Term{graded-module}{分次} $\mathcal O_X[T_1,\ldots,T_d]$-\Term{module}{模}是\Term{perfect-complex}{完美}的。特别地，$\mathcal S$ 作为\Term{graded-module}{分次模}是\Term{finitely-presented}{有限呈示}的。取一个\Term{finite-graded-presentation}{有限分次呈示}
\[
L''\longrightarrow L'\longrightarrow\mathcal S\longrightarrow0.
\]
在 $\mathbb P_X^{d-1}$ 上得到\Term{exact-sequence}{正合列}
\[
\mathcal M''\longrightarrow\mathcal M'\longrightarrow\mathcal O_{X'}\longrightarrow0,
\]
其中 $\mathcal M''=\operatorname{Proj}(L'')$、$\mathcal M'=\operatorname{Proj}(L')$，因此它们在 $\mathbb P_X^{d-1}$ 上是\Term{finite-free}{有限秩自由}的。故对每个 $n$ 都有\Term{commutative-diagram}{交换图}
\[
\begin{tikzcd}[column sep=small]
L''_n\arrow[r]\arrow[d,"\sim"']&L'_n\arrow[r]\arrow[d,"\sim"']&\mathcal J^n\arrow[r]\arrow[d]&0\\
p_*(\mathcal M''(n))\arrow[r]&p_*(\mathcal M'(n))\arrow[r]&p_*(\mathcal O_{X'}(n))\arrow[r]&0,
\end{tikzcd}
\]
其中 $p:\mathbb P_X^{d-1}\to X$ 是\Term{structure-morphism}{结构态射}，$L''_n,L'_n$ 分别是 $L'',L'$ 的 $n$ 次\Term{homogeneous-component}{齐次分量}。上行由假设\Term{exact-sequence}{正合}；前两个竖\Term{arrow}{箭头}是\Term{isomorphism}{同构}，这是\Term{projective-bundle}{射影丛} $P$ 上 $\mathcal O_P(n)$ 的\Term{direct-image}{正像}的经典结论（EGA III 2.1.15）。最后，$n$ 充分大时下行\Term{exact-sequence}{正合}：$\mathcal M'\to\mathcal O_{X'}$ 的\Term{kernel}{核}是\Term{perfect-complex}{完美} $\mathcal O_P$-\Term{module}{模}，可应用前述\Term{lemma}{引理}，得出此\Term{kernel}{核}在与 $\mathcal O_P(n)$ 作\Term{tensor-product}{张量积}后，对于充分大的 $n$ 是 $p_*$-\Term{acyclic-object}{无上同调的}；$\mathcal M''\to\mathcal M'$ 的\Term{kernel}{核}亦然。因此对充分大的 $n$，$\mathcal J^n\to p_*(\mathcal O_{X'}(n))$ 是\Term{isomorphism}{同构}。而 $p_*(\mathcal O_{X'}(n))=f_*(\mathcal O_{X'}(n))$，所以
\SourcePage{454}
已对 $n\geq n_0$ 得到结论。

剩下的工作是\Term{descent}{下降}到 $n=0$。利用\Term{commutative-diagram}{交换图}
\[
\begin{tikzcd}[column sep=small]
0\arrow[r]\arrow[d]&\mathcal J^{n+1}\arrow[r]\arrow[d,"\sim"']&\mathcal J^n\arrow[r]\arrow[d]&\mathcal J^n/\mathcal J^{n+1}\arrow[r]\arrow[d,"\sim"]&0\arrow[d]\\
0\arrow[r]&f_*(\mathcal O_{X'}(n+1))\arrow[r]&f_*(\mathcal O_{X'}(n))\arrow[r]&f_*(\mathcal O_{Y'}(n))\arrow[r]&0
\end{tikzcd}
\]
作\Term{descent}{下降}归纳。下行\Term{exact-sequence}{正合}，因为由已经证明的 i)，对每个 $n\geq0$ 都有 $R^1f_*(\mathcal O_{X'}(n+1))=0$。此外，
\[
f_*(\mathcal O_{Y'}(n))=i_*(g_*(\mathcal O_{Y'}(n)))=i_*(S^n(\mathcal N))=\mathcal J^n/\mathcal J^{n+1}.
\]
所以右\Term{arrow}{箭头}是\Term{isomorphism}{同构}；由归纳假设，左\Term{arrow}{箭头}也是\Term{isomorphism}{同构}，故中间的\Term{arrow}{箭头}亦然。
\end{proof}

\begin{SGAresult}{命题}{3.6}
在 3.1 和 3.3 的假设与记号下，有
\[
f_*\circ f^*=\operatorname{Id}_{K'(X)}.
\]
\end{SGAresult}
\begin{proof}
设 $x\in K'(X)$；由\Term{projection-formula}{投影公式}，
\[
f_*(f^*(x))=f_*(1\cdot f^*(x))=x\cdot f_*(1).
\]
只须证明 $f_*(1)=1$。而 3.5 表明，$\mathbf Rf_*(\mathcal O_{X'})$ 的\Term{cohomology}{上同调}在\Term{degree}{次数} $>0$ 为零，在\Term{degree}{次数} $0$ 等于 $\mathcal O_X$。命题立即得证。
\end{proof}

\begin{SGAresult}{\Term{theorem}{定理}}{3.7}
设 $i:Y\to X$ 是\Term{regular-closed-immersion}{正则闭浸入}，$X'$ 是沿 $Y$ \Term{blow-up}{爆破}所得的\Term{scheme}{概形}，$Y'=X'\times_XY$，从而有\Term{commutative-diagram}{交换图}
\SourcePage{455}
\[
\begin{tikzcd}
Y'\arrow[r,"j"]\arrow[d,"g"']&X'\arrow[d,"f"]\\
Y\arrow[r,"i"']&X.
\end{tikzcd}
\]
假定 $X$ \Term{quasi-compact}{拟紧}且有\Term{ample-sheaf}{丰沛层}。则有\Term{exact-sequence}{正合列}
\[
0\longrightarrow K'(Y)\xrightarrow{u}K'(Y')\times K'(X)\xrightarrow{v}K'(X'),
\]
其中，对 $y\in K'(Y)$ 置
\[
u(y)=(g^*(y)\lambda_{-1}(F),-i_*(y)),
\]
对 $y'\in K'(Y')$、$x\in K'(X)$ 置
\[
v(y',x)=j_*(y')+f^*(x),
\]
记号沿用 3.1 和 3.3。此外，$u$ 有一个\Term{left-inverse}{左逆} $u'$，由 $u'(y',x)=g_*(y')$ 定义。

如果还假定 $X$ \Term{noetherian}{诺特}、$Y$ \Term{regular-scheme}{正则}，则 $v$ \Term{surjective}{满射}，因而有\Term{exact-sequence}{正合列}
\[
0\longrightarrow K'(Y)\xrightarrow{u}K'(Y')\times K'(X)\xrightarrow{v}K'(X')\longrightarrow0.
\]
\end{SGAresult}
\begin{proof}
首先，$u$ 有\Term{left-inverse}{左逆} $u'$，因而尤其是\Term{injective}{单射}。事实上，由\Term{projection-formula}{投影公式}，
\[
g_*(g^*(y)\lambda_{-1}(F))=y\cdot g_*(\lambda_{-1}(F)).
\]
而 $g_*(\lambda_{-1}(F))=1$（第 VI 讲 5.9），于是得到所需结论。

其次，$v\circ u=0$。确实，若 $y\in K'(Y)$，则由 3.4，
\[
v\circ u(y)=v(g^*(y)\lambda_{-1}(F),-i_*(y))
=j_*(g^*(y)\lambda_{-1}(F))-f^*(i_*(y))=0.
\]
所以只须证明 $\operatorname{Ker}(v)\subset\operatorname{Im}(u)$。设 $x\in K'(X)$、$y'\in K'(Y')$ 满足
\SourcePage{456}
\[
j_*(y')+f^*(x)=0.
\]
须找到 $y\in K'(Y)$，使
\[
x=-i_*(y),\qquad y'=g^*(y)\lambda_{-1}(F).\tag{3.7.1}
\]
由于 $u$ 有\Term{left-inverse}{左逆} $u'$，必有
\[
y=u'(y',x)=g_*(y').
\]
因此须证明这个 $y$ 满足 (3.7.1)。

为证明 $x=-i_*(y)$，对关系 $f^*(x)=-j_*(y')$ 施以 $f_*$，按 $y$ 的定义得到
\[
f_*(f^*(x))=-f_*(j_*(y'))=-i_*(g_*(y'))=-i_*(y).
\]
而由命题 3.6，$f_*(f^*(x))=x$，于是得到所求关系。

为完成证明，须证 $y'=g^*(g_*(y'))\lambda_{-1}(F)$。置
\[
z=y'-g^*(g_*(y'))\lambda_{-1}(F).
\]
则
\[
g_*(z)=g_*(y')-g_*(y')g_*(\lambda_{-1}(F))=0.
\]
另一方面，由假设 $j_*(y')=f^*(x)$，故由前面的结果 $j_*(y')=f^*(i_*(y))$；再由命题 3.4，$j_*(y')=j_*(g^*(y)\lambda_{-1}(F))$，所以也有 $j_*(z)=0$。而 $\operatorname{Ker}(j_*)\cap\operatorname{Ker}(g_*)=0$：若 $j_*(z)=0$，则 $j^*(j_*(z))=0$；而 $j^*(j_*(z))=z(1-L)$（命题 2.7）。要使 $z(1-L)=0$，必须存在 $z'\in K'(Y)$，使 $z=z'\lambda_{-1}(F)$（第 VI 讲 5.10）；施以 $g_*$ 便得 $z'=g_*(z)=0$，故 $z=0$。

现在假定 $X$ \Term{noetherian}{诺特}且 $Y$ \Term{regular-scheme}{正则}。那么 $Y'$ 也\Term{regular-scheme}{正则}，因为它是 $Y$ 上的\Term{projective-bundle}{射影丛}；而且对于 $Y'$ 的任一点 $y'$，$X'$ 在 $y'$ 处的\Term{local-ring}{局部环} $\mathcal O_{X',y'}$ 是\Term{regular-ring}{正则环}，因为 $\mathcal O_{Y',y'}$ 是 $\mathcal O_{X',y'}$ 对一个\Term{regular-element}{正则元素}的商（见 EGA $0_{\mathrm{IV}}$ 17.1.8）。
\SourcePage{457}
我们知道，对于有\Term{ample-sheaf}{丰沛层}的\Term{regular-scheme}{正则概形}，用\Term{finite-locally-free}{有限秩局部自由}\Term{sheaf}{层}与用\Term{coherent-sheaf}{凝聚层}构造的\Term{grothendieck-group}{Grothendieck 群}相同。$Y'$ 尤其满足这一条件。

设 $\mathcal E$ 是\Term{support}{支集}包含在 $Y'$ 中的\Term{coherent-sheaf}{凝聚} $\mathcal O_{X'}$-\Term{module}{模}；则 $\mathcal E$ 是\Term{perfect-complex}{完美} $\mathcal O_{X'}$-\Term{module}{模}。事实上，在 $Y'$ 的点处如此，因为 $X'$ 在这些点\Term{regular-scheme}{正则}；在 $X'-Y'$ 的点处也如此，因为在那里它为零。

若 $E$ 为 $\mathcal E$ 在 $K'(X')$ 中的类，则 $E$ 属于 $j_*$ 的像。确实，可用各 $\mathcal J'^n\mathcal E$ 为 $\mathcal E$ 作\Term{filtration}{滤过}。此\Term{filtration}{滤过}\Term{finite}{有限}，因为 $X'$ \Term{noetherian}{诺特}，所以存在整数 $n_0$，使 $\mathcal J'^{n_0}\mathcal E=0$；此\Term{filtration}{滤过}的各商都是 $\mathcal O_{Y'}$-\Term{module}{模}，故其在 $K'(X')$ 中的类属于 $j_*$ 的像。由于 $E$ 是这些商的类之和，$E$ 也属于该像。

现设 $x\in K'(X')$，取类为 $x$ 的\Term{perfect-complex}{完美复形} $\mathcal M^\bullet$。考虑\Term{distinguished-triangle}{区别三角}
\[
\begin{tikzcd}
&\mathcal N^\bullet\arrow[dl]\arrow[dr]&\\
\mathcal M^\bullet\arrow[rr]&&\mathbf Lf^*(\mathbf Rf_*(\mathcal M^\bullet)),
\end{tikzcd}
\]
其中水平\Term{arrow}{箭头}是\Term{canonical-map}{典范同态}。此\Term{arrow}{箭头}在 $Y'$ 的补集上是\Term{isomorphism}{同构}，故 $\mathcal N^\bullet$ 的\Term{cohomology-sheaf}{上同调层}是\Term{support}{支集}包含在 $Y'$ 中的\Term{coherent-sheaf}{凝聚层}。由前面的结果，它们都是\Term{perfect-complex}{完美} $\mathcal O_{X'}$-\Term{module}{模}，且其类属于 $j_*$ 的像。因此
\[
x-f^*(f_*(x))\in\operatorname{Im}(j_*),
\]
从而\Term{homomorphism}{同态} $v$ \Term{surjective}{满射}。
\end{proof}
\SourcePage{458}
\textbf{注记 3.8。} 在 $X$ \Term{noetherian}{诺特}、$Y$ \Term{regular-scheme}{正则}的情形下，3.7 完全用 $\lambda$-\Term{lambda-ring}{环} $K'(X)$、$K'(Y)$，这些\Term{ring}{环}之间的映射 $i_*$、$i^*$，以及 $K'(Y)$ 的元素 $N$，确定了\Term{abelian-group}{阿贝尔群} $K'(X')$。事实上，由第 VI 讲 1.1，$K'(Y')$ 可以用 $K'(Y)$ 和 $N$ 算出。下一节将看到，$K'(X')$ 的\Term{multiplication}{乘法}结构也由这些数据确定；若与 $\mathbb Q$ 作\Term{tensor-product}{张量积}，其 $\lambda$-结构亦然（这个\Term{restriction}{限制}很可能没有必要：见 4.5）。

\SGAsection{4}{正则浸入与 $K'$ 的滤过}
本节保持 3.1 的记号和假设。
\begin{SGAresult}{\Term{lemma}{引理}}{4.1}
假定 $Y$ 上存在一个\Term{finite-locally-free}{有限秩局部自由}\Term{sheaf}{层} $\mathcal N'$，使得若 $N,N'$ 分别为 $\mathcal N,\mathcal N'$ 在 $K'(Y)$ 中的类，则 $N=N'+2$。那么有如下同余：
\[
\lambda_{-1}(F)\equiv0\pmod{(1-L)}.
\]
\end{SGAresult}
\begin{proof}
设 $d$ 为 $\mathcal N$ 的\Term{rank}{秩}，可假定它为常数；那么 $\mathcal F$ 的\Term{rank}{秩}为 $d-1$。由定义，
\[
\begin{aligned}
\lambda_{-1}(F)&=\sum_{i\geq0}(-1)^i\lambda^i(F)\\
&=(-1)^{d-1}\lambda^{d-1}(F-1)\\
&=(-1)^{d-1}\lambda^{d-1}(F+L-1-L)\\
&=(-1)^{d-1}\lambda^{d-1}(N-2+1-L).
\end{aligned}
\]
而 $\lambda_t(1-L)=\lambda_t(1)/\lambda_t(L)=(1+t)/(1+Lt)$，所以对于每个 $i\geq1$，$\lambda^i(1-L)$ 都可被 $1-L$ 整除。因此
\[
\begin{aligned}
\lambda_{-1}(F)&\equiv(-1)^{d-1}\lambda^{d-1}(N-2)&&\pmod{(1-L)}\\
&\equiv(-1)^{d-1}\lambda^{d-1}(N')&&\pmod{(1-L)}.
\end{aligned}
\]
\SourcePage{459}
$\mathcal N'$ 的\Term{rank}{秩}为 $d-2$，故 $\lambda^{d-1}(N')=0$，结论得证。
\end{proof}

\textbf{4.2.} 设 $i:Y\to X$ 是\Term{regular-immersion}{正则浸入}，$r\geq1$ 为整数，$s:X\to\mathbb P_X^r$ 是由\Term{homomorphism}{同态} $(\mathcal O_X)^{r+1}\to\mathcal O_X$ 定义的\Term{section}{截面}，其中此\Term{homomorphism}{同态}取对后 $r$ 个坐标生成的\Term{subsheaf}{子层}的商。若 $\mathbb P_X^r$ 由\Term{graded-algebra}{分次} $\mathcal O_X$-\Term{algebra}{代数} $\mathcal O_X[T_0,\ldots,T_r]$ 定义，则定义\Term{section}{截面} $s$ 的\Term{ideal}{理想} $\mathcal I$ 由\Term{section}{截面} $T_1,\ldots,T_r$ 生成；因而 $\mathcal I/\mathcal I^2$ \Term{isomorphism}{同构}于 $(\mathcal O_X[T_0])^r$，所以相应的 $\mathcal O_P$ 的\Term{ideal}{理想}\Term{isomorphism}{同构}于 $(\mathcal O_X)^r$。

显然\Term{section}{截面} $s$ 是\Term{regular-immersion}{正则浸入}，故 $s\circ i$ 亦然，并可写出相应的\Term{conormal-sheaf}{余法层}\Term{exact-sequence}{正合列}（1.7）。于是以 $\mathcal N$ 表示 $i$ 的\Term{conormal-sheaf}{余法层}，$\mathcal N_1$ 表示 $s\circ i$ 的\Term{conormal-sheaf}{余法层}，在 $K'(Y)$ 中得到等式
\[
N_1=N+r.
\]
所以对于任一\Term{regular-immersion}{正则浸入} $i$，都可将它与一个\Term{section}{截面} $i:X\to\mathbb P_X^2$，或 $X\to\mathbb P_{\mathbb P_X^1}^1$ 复合，得到满足\Term{lemma}{引理} 4.1 所述性质的\Term{regular-immersion}{正则浸入}。

\begin{SGAresult}{\Term{theorem}{定理}}{4.3}
设 $i:Y\to X$ 为\Term{regular-closed-immersion}{正则闭浸入}，$N$ 为 $Y$ 在 $X$ 中的\Term{conormal-sheaf}{余法层}在 $K'(Y)$ 中的类。假定 $X$ \Term{quasi-compact}{拟紧}且有\Term{ample-sheaf}{丰沛层}。则沿用 3.3 的记号，对每个 $y\in K'(Y)$ 及每个 $p\geq1$，有
\[
i_*(\lambda^p(N,y))\equiv\lambda^p(i_*(y))\pmod{\text{挠}}.\tag{4.3.1}
\]
其中 $\lambda^p(N,y)$ 是第 V 讲 5.3 定义的元素。
\end{SGAresult}
\begin{proof}
i) 首先假定满足 4.1 的假设；此时将证明 $i_*(\lambda^p(N,y))=\lambda^p(i_*(y))$。由于 $f_*\circ f^*=\operatorname{Id}$，这个关系\Term{equivalence}{等价}于
\SourcePage{460}
\[
f^*(i_*(\lambda^p(N,y)))=f^*(\lambda^p(i_*(y))).
\]
一方面有
\[
\begin{aligned}
f^*(\lambda^p(i_*(y)))&=\lambda^p(f^*(i_*(y)))\\
&=\lambda^p(j_*(g^*(y)\lambda_{-1}(F)))\qquad\text{（3.4）}.
\end{aligned}
\]
另一方面，
\[
f^*(i_*(\lambda^p(N,y)))=j_*(g^*(\lambda^p(N,y))\lambda_{-1}(F))\qquad\text{（3.4）}.
\]
由于 $g^*(N)=F+L$，可应用公式第 V 讲 (5.4.2)：
\[
g^*(\lambda^p(N,y))\lambda_{-1}(F)
=\lambda^p(g^*(N),g^*(y))\lambda_{-1}(F)
=\lambda^p(L,g^*(y)\lambda_{-1}(F)).
\]
置 $z=g^*(y)\lambda_{-1}(F)$，问题就归结为证明
\[
j_*(\lambda^p(L,z))=\lambda^p(j_*(z)),
\]
即对\Term{immersion}{浸入} $j$ 和元素 $z$ 的相应关系。由假设，$\lambda_{-1}(F)$ 是 $1-L$ 的倍数；可取 $z'$ 使 $z=z'(1-L)$。于是由 2.7，$z=j^*(u')$，其中 $u'=j_*(z')$；由于 $L=j^*(\mathcal J')$，最终只须证明
\[
j_*(j^*(\lambda^p(\mathcal J',u')))=\lambda^p(j_*(j^*(u'))).
\]
\Term{exact-sequence}{正合列}
\[
0\longrightarrow\mathcal J'\longrightarrow\mathcal O_{X'}\longrightarrow\mathcal O_{Y'}\longrightarrow0
\]
给出关系
\[
j_*(1)=1-\mathcal J'.\tag{4.3.2}
\]
利用\Term{projection-formula}{投影公式}，待证关系遂可写为
\SourcePage{461}
\[
\lambda^p(\mathcal J',u')(1-\mathcal J')=\lambda^p(u'(1-\mathcal J')).
\]
由于 $1-\mathcal J'=\lambda_{-1}(\mathcal J')$，这正是定义 $\lambda^p(\mathcal J',u')$ 所用的公式。

ii) 现考虑\Term{generic}{一般}情形。先将 $i$ 与一个\Term{section}{截面} $X\to\mathbb P_X^1$ 复合，再与一个\Term{section}{截面} $\mathbb P_X^1\to\mathbb P_{\mathbb P_X^1}^1$ 复合，使复合的\Term{immersion}{浸入}满足 4.1 的条件（见 4.2）。由前面的证明，命题 4.3 对此复合\Term{immersion}{浸入}成立；所以只须将其“降回”到 $X$。先降到 $\mathbb P_X^1$，再降到 $X$，因此只须应用两次下述\Term{lemma}{引理}（由 J. P. Jouanolou 提出）。

\begin{SGAresult}{\Term{lemma}{引理}}{}
在 4.3 的条件下，设 $s:X\to\mathbb P_X^1$ 是 4.2 定义的\Term{section}{截面}。若命题 4.3 对 $s\circ i$ 成立，则对 $i$ 也成立。
\end{SGAresult}
假定 4.3 对 $s\circ i$ 成立，则对于每个 $p\geq1$ 和 $y\in K'(Y)$，都有
\[
(s\circ i)_*(\lambda^p(N+1,y))\equiv\lambda^p((s\circ i)_*(y))\pmod{\text{挠}}.\tag{4.3.3}
\]
\Term{section}{截面} $s$ 由一个\Term{isomorphism}{同构}于 $\mathcal O_P(-1)$ 的\Term{ideal}{理想}定义（4.2）；以 $\xi$ 表示它在 $K'(\mathbb P_X^1)$ 中的类。则对每个 $x\in K'(X)$，都有
\[
s_*(x)=p^*(x)(1-\xi),
\]
其中 $p:\mathbb P_X^1\to X$ 为\Term{structure-morphism}{结构态射}。事实上，由 $p\circ s=\operatorname{Id}_X$ 得 $x=s^*(p^*(x))$，所以
\[
s_*(x)=s_*(s^*(p^*(x)))=p^*(x)s_*(1);
\]
而 $s_*(1)=1-\xi$，遂得所求关系。关系 (4.3.3) 于是写成
\[
\begin{aligned}
(s\circ i)_*(\lambda^p(N+1,y))
&\equiv\lambda^p(p^*(i_*(y))(1-\xi))&&\pmod{\text{挠}}\\
&\equiv\lambda^p(\xi,p^*(i_*(y)))(1-\xi)&&\pmod{\text{挠}}.
\end{aligned}
\]
\SourcePage{462}
而 $\lambda^p(\xi,p^*(i_*(y)))-\lambda^p(1,p^*(i_*(y)))$ 可被 $1-\xi$ 整除。为看出这一点，可在定义 $\lambda^p(\xi,x)$ 所用的泛 $\lambda$-\Term{lambda-ring}{环}中工作（第 V 讲第 5 节）；此时 $\lambda^p(\xi,x)-\lambda^p(1,x)$ 是关于 $\xi$ 的\Term{polynomial}{多项式}，在 $\xi=1$ 时为零，故可被 $\xi-1$ 整除。另一方面，$\xi$ 在 $K'(X)$ 上的依赖方程为 $(1-\xi)^2=0$（第 VI 讲 (1.13.1)）。于是
\[
\begin{aligned}
(s\circ i)_*(\lambda^p(N+1,y))
&\equiv\lambda^p(1,p^*(i_*(y)))(1-\xi)&&\pmod{\text{挠}}\\
&\equiv p^*(\lambda^p(1,i_*(y)))(1-\xi)&&\pmod{\text{挠}}.
\end{aligned}
\]
施以 $p_*$，因 $p_*(\xi)=0$，得到
\[
i_*(\lambda^p(N+1,y))\equiv\lambda^p(1,i_*(y))\pmod{\text{挠}}.\tag{4.3.4}
\]
将第 V 讲 5.6 的公式用于 $L=1$，可知 $\lambda^p(N+1,y)$ 是级数
\[
\frac{\displaystyle\sum_{k\geq1}\lambda^k(N,y)kt^k}
{\displaystyle1+\sum_{k\geq1}\lambda^k(N,y)\lambda_{-1}(N)t^k}
\]
中 $t^p$ 的系数。同样，将该公式用于 $N=0$、$L=1$，可知 $\lambda^p(1,y)$ 是级数
\[
\frac{\displaystyle\sum_{k\geq1}\lambda^k(y)kt^k}
{\displaystyle\sum_{k\geq0}\lambda^k(y)t^k}
\]
中 $t^p$ 的系数。所以同余式 (4.3.4) 可写为
\[
i_*\left(\frac{\displaystyle\sum_{k\geq1}k\lambda^k(N,y)t^k}
{\displaystyle1+\sum_{k\geq1}\lambda^k(N,y)\lambda_{-1}(N)t^k}\right)
\equiv\frac{\displaystyle\sum_{k\geq1}k\lambda^k(i_*(y))t^k}
{\displaystyle\sum_{k\geq0}\lambda^k(i_*(y))t^k}\pmod{\text{挠}}.
\]
\SourcePage{463}
应用公式 2.8，得到
\[
\frac{\displaystyle\sum_{k\geq0}ki_*(\lambda^k(N,y))t^k}
{\displaystyle1+\sum_{k\geq1}i_*(\lambda^k(N,y))t^k}
\equiv\frac{\displaystyle\sum_{k\geq0}k\lambda^k(i_*(y))t^k}
{\displaystyle\sum_{k\geq0}\lambda^k(i_*(y))t^k}\pmod{\text{挠}}.
\]
若置
\[
u(t)=1+\sum_{k\geq1}i_*(\lambda^k(N,y))t^k,\qquad
v(t)=\sum_{k\geq0}\lambda^k(i_*(y))t^k,
\]
则最后一个关系还可写成
\[
\frac{d}{dt}\left(\frac{u(t)}{v(t)}\right)\equiv0\pmod{\text{挠}}.
\]
由于 $u(t)/v(t)$ 的常数项为 $1$，且我们所写的是模\Term{torsion}{挠}的同余式，故可积分而得同余式 $u(t)\equiv v(t)$。这就给出了所述各同余式。
\end{proof}

\begin{SGAresult}{推论}{4.4}
在 4.3 的记号与假设下，对每个 $p\geq1$ 和每个 $y\in K'(Y)$，都有
\[
i_*(\gamma^p(N,y))\equiv\gamma^p(i_*(y))\pmod{\text{挠}}.
\]
\end{SGAresult}
此推论由 4.3 立即得到，只须利用各 $\gamma^p$ \Term{representable}{可表}成各 $\lambda^p$ 的整系数\Term{linear-combination}{线性组合}。

\textbf{注记 4.5。} 4.3 的证明本身使人相信，(4.3.1) 应当作为等式成立，而不仅是模\Term{torsion}{挠}成立。在 [RRR] 中，借助 $\lambda^p(N,x)$ 的几何刻画，已在\Term{characteristic}{特征} $0$ 下证明了此等式。如果能够在\Term{perfect-complex}{完美复形}上定义\Term{exterior-power}{外幂}运算，这种方法很可能可以推广到任意\Term{characteristic}{特征}（参见第 XIV 讲）。

\SourcePage{464}
\begin{SGAresult}{\Term{theorem}{定理}}{4.6}
设 $X$ 是有\Term{ample-sheaf}{丰沛层}的\Term{quasi-compact}{拟紧}\Term{scheme}{概形}，$d>0$ 为整数，$i:Y\to X$ 是\Term{codimension}{余维数}为 $d$ 的\Term{regular-immersion}{正则浸入}。则对每个 $k$ 都有
\[
i_*\bigl((\operatorname{Fil}^kK'(Y))\otimes_{\mathbb Z}\mathbb Q\bigr)
\subset(\operatorname{Fil}^{k+d}K'(X))\otimes_{\mathbb Z}\mathbb Q.
\]
\end{SGAresult}
\begin{proof}
首先注意 $i_*(K'(Y))\subset\operatorname{Fil}^1K'(X)$。事实上，$X-Y$ 在 $X$ 中\Term{dense}{稠密}：否则 $Y$ 会包含一个\Term{affine-open}{仿射开集}，在这个\Term{open-subset}{开集}上，定义 $Y$ 的\Term{ideal}{理想} $\mathcal J$ 将是一个\Term{nil-ideal}{幂零元理想}；又因其\Term{finite-type}{有限生成}，故是\Term{local-scheme}{局部}\Term{nilpotent-ideal}{幂零理想}。但这不可能，因为各 $\mathcal J^n/\mathcal J^{n+1}$ 是\Term{rank}{秩} $\ne0$ 的\Term{finite-locally-free}{有限秩局部自由} $\mathcal O_Y$-\Term{module}{模}。$i_*(K'(Y))$ 的元素在 $K'(X)$ 中的\Term{augmentation}{增广}是 $X$ 上的一个\Term{locally-constant-function}{局部常值函数}；它在 $X-Y$ 上为零，故在整个 $X$ 上为零。

对任意\Term{abelian-group}{阿贝尔群} $A$，记 $A_{\mathbb Q}=A\otimes_{\mathbb Z}\mathbb Q$。设 $x\in\operatorname{Fil}^kK'(Y)_{\mathbb Q}$；可将 $x$ 写为
\[
x=\sum a\gamma^{i_1}(x_1)\cdots\gamma^{i_n}(x_n),
\]
其中每个 $i$ 都有 $\varepsilon(x_i)=0$，$a\in\mathbb Q$，且 $i_1+\cdots+i_n\geq k$。各 $x_i$ 都\Term{representable}{可表}为\Term{finite}{有限}多个\Term{finite-locally-free}{有限秩局部自由} $\mathcal O_Y$-\Term{module}{模} $\mathcal E_\alpha$ 的类 $E_\alpha$ 的整系数\Term{linear-combination}{线性组合}。设 $R$ 为 $K'(Y)_{\mathbb Q}$ 中由这些 $E_\alpha$ 及 $N$ 生成的 $\mathbb Q$-$\lambda$-\Term{subalgebra}{子代数}。于是
\[
x\in\operatorname{Fil}^kR.
\]
此外，$R$ 的 $\lambda$-\Term{lambda-ring}{环}\Term{filtration}{滤过}是\Term{discrete-filtration}{离散}的（第 VI 讲 6.2）；取 $m$ 使 $\operatorname{Fil}^mR=0$。最后，把 $X$ \Term{resolution}{分解}为\Term{finite}{有限}个不交\Term{open-subset}{开集}之和，并在每个\Term{open-subset}{开集}上分别证明（$K'$ 与\Term{finite}{有限}\Term{direct-sum}{直和}可交换），可假定各 $\mathcal E_\alpha$ 有常\Term{rank}{秩} $r_\alpha$。此时 $R$ 的\Term{augmentation-ring}{增广环}就是 $\mathbb Q$。
\SourcePage{465}
将证明对每个 $i$ 都有
\[
i_*(\operatorname{Fil}^iR)\subset\operatorname{Fil}^{i+d}(K'(X)_{\mathbb Q}).
\]
这尤其给出 $i_*(x)\in\operatorname{Fil}^{k+d}K'(X)$，因而证明\Term{theorem}{定理}。

对 $i$ 作\Term{descent}{下降}归纳；$i=m$ 的情形显然。设 $y\in\operatorname{Fil}^iR$，对每个 $j$ 置
\[
a_j=(-1)^{j-1}(j-1)!.
\]
由第 V 讲 6.10，
\[
\gamma^{d+i}(N,y)-a_{d+i}y\in\operatorname{Fil}^{i+1}R.
\]
应用归纳假设
\[
i_*(\operatorname{Fil}^{i+1}R)\subset\operatorname{Fil}^{i+d+1}K'(X)_{\mathbb Q},
\]
遂有
\[
i_*(\gamma^{d+i}(N,y)-a_{d+i}y)\in\operatorname{Fil}^{i+d+1}K'(X)_{\mathbb Q}.
\]
由 4.4，得到
\[
\gamma^{d+i}(i_*(y))-a_{d+i}i_*(y)\in\operatorname{Fil}^{i+d+1}K'(X)_{\mathbb Q}.
\]
由于 $i_*(y)$ 的\Term{augmentation}{增广}为零，$\gamma^{d+i}(i_*(y))\in\operatorname{Fil}^{i+d}K'(X)$。在模\Term{torsion}{挠}意义下计算，便得到
\[
i_*(y)\in\operatorname{Fil}^{i+d}K'(X)_{\mathbb Q}.
\]
\end{proof}
\textbf{注记 4.6.1。} 可见，即使 4.3 在不与 $\mathbb Q$ 作\Term{tensor-product}{张量积}时成立，上述证明也只在模\Term{torsion}{挠}意义下适用。例如可参看第 XIV 讲 4.6。

\SourcePage{466}
\textbf{4.7.} 回到 3.7 的情形，那里有\Term{commutative-diagram}{交换图}
\[
\begin{tikzcd}
Y'\arrow[r,"j"]\arrow[d,"g"']&X'\arrow[d,"f"]\\
Y\arrow[r,"i"']&X,
\end{tikzcd}
\]
其中 $i$ 是\Term{regular-immersion}{正则浸入}，$X'$ 是沿 $Y$ \Term{blow-up}{爆破}所得的 $X$-\Term{scheme}{概形}，$Y'=X'\times_XY$。假定 $X$ \Term{noetherian}{诺特}且有\Term{ample-sheaf}{丰沛层}，$Y$ \Term{regular-scheme}{正则}。由 3.7 有\Term{exact-sequence}{正合列}
\[
0\longrightarrow K'(Y)\xrightarrow{u}K'(Y')\times K'(X)\xrightarrow{v}K'(X')\longrightarrow0.
\]
现在将说明，这个\Term{exact-sequence}{正合列}使我们能够用 $\lambda$-\Term{lambda-ring}{环} $K'(X)$、$K'(Y)$，映射 $i^*$、$i_*$ 及 $K'(Y)$ 的元素 $N$，描述由 $X,Y,X',Y'$ 的 $K'$ 构成的整个情形（目前至少在模\Term{torsion}{挠}意义下如此）。

因此设 $K,K'$ 为两个 $\lambda$-\Term{lambda-ring}{环}，$i_*:K\to K'$ 为\Term{additive-homomorphism}{加法同态}，$i^*:K'\to K$ 为 $\lambda$-\Term{lambda-homomorphism}{同态}，$d$ 为整数，$N\in K$ 满足 $i>d$ 时 $\lambda^i(N)=0$，且 $\lambda^d(N)$ \Term{invertible-element}{可逆}。假定下列性质成立：

 i) $i_*$ 是 $K'$-\Term{linear-map}{线性的}（“\Term{projection-formula}{投影公式}”）；

 ii) $i^*\circ i_*$ 是“乘以 $\lambda_{-1}(N)$”的\Term{homomorphism}{同态}；

 iii) 对每个 $n\geq1$ 及每个 $y\in K$，$i_*(\lambda^n(N,y))=\lambda^n(i_*(y))$（关系 (4.3.1)）。

先利用第 VI 讲 1.1 构造一个 $K$-$\lambda$-\Term{algebra}{代数} $K_1$，它扮演 $K'(Y')$ 的角色。设 $L$ 为\Term{indeterminate}{不定元}，$K[L]$ 为由\Term{generator}{生成元} $L$ 在关系 $i\geq2$ 时 $\lambda^i(L)=0$ 下生成的 $K$-$\lambda$-\Term{algebra}{代数}（第 V 讲 4.9）。置
\[
P(L)=\lambda^d(N-L)=\sum_{i=0}^d(-1)^i\lambda^{d-i}(N)L^i.
\]
\SourcePage{467}
由 $P$ 生成的\Term{ideal}{理想}是一个 $\lambda$-\Term{lambda-ideal}{理想}。为看出这一点，先注意 $i>d$ 时 $\lambda^i(N)=0$ 蕴涵 $j\geq d$ 时各 $\lambda^j(N-L)$ 都是 $\lambda^d(N-L)$ 的倍数。在 $1+\widehat{(K/(P))[[t]]}^{+}$ 中，$\lambda_t(N-L)$ 的像是\Term{degree}{次数} $<d$ 的\Term{polynomial}{多项式}；一个泛论证表明，它在 $i\geq d$ 时的各 $\lambda^i$ 均等于 $1$（参见第 V 讲 4.8）。由于在 $1+\widehat{K[[t]]}^{+}$ 中有
\[
\lambda_t(\lambda^d(N-L))=\lambda^d(\lambda_t(N-L)),
\]
可知对 $i\geq1$，各 $\lambda^i(\lambda^d(N-L))$ 都是 $P$ 的倍数。由第 V 讲 2.4.2 的关系，这意味着由 $P$ 生成的\Term{ideal}{理想}是 $\lambda$-\Term{lambda-ideal}{理想}。

于是置 $K_1=K/(P)$，以 $g^*$ 表示自然\Term{homomorphism}{同态} $K\to K_1$；它是 $\lambda$-\Term{lambda-homomorphism}{同态}。另一方面，$K_1$ 是以 $1,L,\ldots,L^{d-1}$ 为\Term{basis}{基}的\Term{free-module}{自由} $K$-\Term{module}{模}。由关系 $\lambda^d(N)=\lambda^{d-1}(N-L)L$ 以及 $\lambda^d(N)$ \Term{invertible-element}{可逆}的假设，$L$ \Term{invertible-element}{可逆}。所以 $1,L^{-1},\ldots,L^{-d+1}$ 也是 $K_1$ 在 $K$ 上的一组\Term{basis}{基}，由此定义映射 $g_*:K_1\to K$，将每个元素
\[
y'=y_0+y_1L^{-1}+\cdots+y_{d-1}L^{-d+1}
\]
映到 $K$ 的元素 $y_0$。显然，如此定义的映射是 $K$-\Term{linear-map}{线性}的\Term{additive-homomorphism}{加法同态}，且 $g_*\circ g^*=\operatorname{Id}$。

接着在 $K_1\times K'$ 上定义一个 $\lambda$-\Term{lambda-ring}{环}结构，置
\[
(y'_1,x_1)*(y'_2,x_2)
=(y'_1y'_2(1-L)+x_1y'_2+x_2y'_1,\ x_1x_2),\tag{4.7.1}
\]
\[
\lambda_*^n(y',x)=\left(\sum_{p=1}^n\lambda^{n-p}(x)\lambda^p(L,y'),\ \lambda^n(x)\right).\tag{4.7.2}
\]
在这些公式中，通过\Term{homomorphism}{同态} $g^*\circ i^*$ 将 $K_1$ 看作 $K'$-\Term{algebra}{代数}。再对 $y_1,y_2\in K$ 及 $n\geq1$ 置
\[
y_1*y_2=-y_1y_2\lambda_{-1}(N),\tag{4.7.3}
\]
\[
\lambda_*^n(y_1)=-\lambda^n(N,-y_1),\tag{4.7.4}
\]
\SourcePage{468}
并定义\Term{homomorphism}{同态} $u:K\to K_1\times K'$，令
\[
u(y)=(g^*(y)\lambda_{-1}(F),-i_*(y)),
\]
其中 $F=N-L$。

\begin{SGAresult}{命题}{4.7.5}
由 (4.7.1) 与 (4.7.2) 定义的运算在 $K_1\times K'$ 上给出一个 $\lambda$-\Term{lambda-ring}{环}结构；$u$ 是\Term{injective}{单射}\Term{homomorphism}{同态}，将带有 (4.7.3)、(4.7.4) 所定义的运算的 $K$ 识别为 $K_1\times K'$ 的一个 $\lambda$-\Term{lambda-ideal}{理想}。
\end{SGAresult}
$K_1\times K'$ 为 $\lambda$-\Term{lambda-ring}{环}是直接可证的，例如可参看第 V 讲 5.5，那里给出了完全类似的构造。命题第二部分的证明也没有困难：利用前面的关系、由假设立即得到的关系
\[
i_*(y_1y_2\lambda_{-1}(N))=i_*(y_1)i_*(y_2),
\]
关系 $g_*(\lambda_{-1}(F))=1$（见第 VI 讲 5.9）以及第 V 讲第 5 节的各关系即可。

于是置
\[
K'_1=(K_1\times K')/K,
\]
以 $v$ 表示\Term{canonical-map}{典范同态} $K_1\times K'\to K'_1$。通过\Term{quotient-structure}{商结构}，将 $K'_1$ 看作 $\lambda$-\Term{lambda-ring}{环}。

再定义各\Term{additive-homomorphism}{加法同态}
\[
\begin{aligned}
j_* &:K_1\longrightarrow K'_1,&j^* &:K'_1\longrightarrow K_1,\\
f_* &:K'_1\longrightarrow K',&f^* &:K'\longrightarrow K'_1.
\end{aligned}
\]
\Term{homomorphism}{同态} $j_*$ 是复合 $K_1\to K_1\times K'\to K'_1$。同样，\Term{homomorphism}{同态} $f^*$ 是复合 $K'\to K_1\times K'\to K'_1$。于是
\[
v=j_*+f^*.
\]
\SourcePage{469}
再定义\Term{homomorphism}{同态} $K_1\times K'\to K_1$，将 $(y',x)$ 映到 $y'(1-L)+g^*\circ i^*(x)$。这个\Term{homomorphism}{同态}在 $K$ 上为零，故定义出一个\Term{homomorphism}{同态} $j^*:K'_1\to K_1$。最后定义\Term{homomorphism}{同态} $K_1\times K'\to K'$，将 $(y',x)$ 映到 $i_*\circ g_*(y')+x$；这个\Term{homomorphism}{同态}在 $K$ 上为零，故定义出一个\Term{homomorphism}{同态} $f_*:K'_1\to K'$。

\begin{SGAresult}{命题}{4.7.6}
 i) $j^*$ 与 $f^*$ 是 $\lambda$-\Term{lambda-homomorphism}{同态}；

 ii) $j_*$ 是 $K'_1$-\Term{linear-map}{线性的}，$f_*$ 是 $K'$-\Term{linear-map}{线性的}；

 iii) $j^*\circ j_*$ 是“乘以 $1-L$”的\Term{homomorphism}{同态}；

 iv) 对每个 $y'\in K_1$ 和每个 $p\geq1$，$j_*(\lambda^p(L,y'))=\lambda^p(j_*(y'))$；

 v) $f_*\circ f^*=\operatorname{Id}$。
\end{SGAresult}
只须按上述定义逐项运算即可。

回到上面回顾的几何情形，可见 $K'(X')$ 可由 $K'(Y)$、$K'(X)$，映射 $i^*$、$i_*$ 以及 $K'(Y)$ 的元素 $N$，按上述方法构造出来。不过，只要关系 (4.3.1) 仍只在模\Term{torsion}{挠}意义下得到证明，$K'(X')$ 上的 $\lambda$ 运算就也只在相差\Term{torsion}{挠}的意义下与上面在 $K'_1$ 上定义的 $\lambda$ 运算一致。

\begin{SGAresult}{\Term{theorem}{定理}}{4.8}
设 $i:Y\to X$ 是常\Term{codimension}{余维数}为 $d$ 的\Term{regular-closed-immersion}{正则闭浸入}，$X'$ 是沿 $Y$ \Term{blow-up}{爆破}所得的\Term{scheme}{概形}，$Y'=X'\times_XY$，从而有\Term{commutative-diagram}{交换图}
\[
\begin{tikzcd}
Y'\arrow[r,"j"]\arrow[d,"g"']&X'\arrow[d,"f"]\\
Y\arrow[r,"i"']&X.
\end{tikzcd}
\]
假定 $X$ \Term{quasi-compact}{拟紧}且有\Term{ample-sheaf}{丰沛层}。则对每个 $k$，有\Term{exact-sequence}{正合列}
\SourcePage{470}
\[
0\longrightarrow\operatorname{Gr}^{k-d}K'(Y)_{\mathbb Q}
\xrightarrow{u^{\mathrm{gr}}}\operatorname{Gr}^{k-1}K'(Y')_{\mathbb Q}\times\operatorname{Gr}^kK'(X)_{\mathbb Q}
\xrightarrow{v^{\mathrm{gr}}}\operatorname{Gr}^kK'(X')_{\mathbb Q},
\]
其中，对 $y\in\operatorname{Gr}^{k-d}K'(Y)_{\mathbb Q}$ 置
\[
u^{\mathrm{gr}}(y)=(g^{\mathrm{gr}}(y)(-1)^{d-1}c^{d-1}(F),-i_{\mathrm{gr}}(y)),
\]
对 $y'\in\operatorname{Gr}^{k-1}K'(Y')_{\mathbb Q}$ 和 $x\in\operatorname{Gr}^kK'(X)_{\mathbb Q}$ 置
\[
v^{\mathrm{gr}}(y',x)=j_{\mathrm{gr}}(y')+f^{\mathrm{gr}}(x).
\]
记号沿用 3.1 和 3.3，$c^{d-1}(F)$ 表示 $F$ 的第 $(d-1)$ 个\Term{chern-class}{Chern 类}（第 V 讲 6.7）。此外，$u^{\mathrm{gr}}$ 有一个\Term{left-inverse}{左逆} $u'^{\mathrm{gr}}$，由 $u'^{\mathrm{gr}}(y',x)=g_{\mathrm{gr}}(y')$ 定义。

如果还假定 $X$ \Term{noetherian}{诺特}、$Y$ \Term{regular-scheme}{正则}，则有\Term{exact-sequence}{正合列}
\[
0\longrightarrow\operatorname{Gr}^{k-d}K'(Y)_{\mathbb Q}
\xrightarrow{u^{\mathrm{gr}}}\operatorname{Gr}^{k-1}K'(Y')_{\mathbb Q}\times\operatorname{Gr}^kK'(X)_{\mathbb Q}
\xrightarrow{v^{\mathrm{gr}}}\operatorname{Gr}^kK'(X')_{\mathbb Q}\longrightarrow0.
\]
\end{SGAresult}
\begin{proof}
注意，在 $Y$ \Term{regular-scheme}{正则}的情形下，可以利用 4.7 对此情形的描述，用 $\lambda$-\Term{lambda-ring}{环}理论给出一个纯\Term{algebra}{代数}的证明。这里将仿照 3.7 的证明。

我们知道，“\Term{inverse-image}{逆像}”\Term{homomorphism}{同态} $f^*,g^*$ 是 $\lambda$-\Term{lambda-homomorphism}{同态}，因而与 $\lambda$-\Term{lambda-ring}{环}\Term{filtration}{滤过}相容，所以在\Term{associated-graded}{伴随分次}上定义\Term{homomorphism}{同态}，即
\[
\begin{aligned}
f^{\mathrm{gr}}&:\operatorname{Gr}^kK'(X)\longrightarrow\operatorname{Gr}^kK'(X'),\\
g^{\mathrm{gr}}&:\operatorname{Gr}^{k-d}K'(Y)\longrightarrow\operatorname{Gr}^{k-d}K'(Y'),
\end{aligned}
\]
以及与 $\mathbb Q$ 作\Term{tensor-product}{张量积}后所得的相应\Term{homomorphism}{同态}。另一方面，4.6 表明，与 $\mathbb Q$ 作\Term{tensor-product}{张量积}后，\Term{homomorphism}{同态} $i_*,j_*$ 可\Term{descent}{下降}到\Term{associated-graded}{伴随分次}，定义
\[
\begin{aligned}
i_{\mathrm{gr}}&:\operatorname{Gr}^{k-d}K'(Y)_{\mathbb Q}\longrightarrow\operatorname{Gr}^kK'(X)_{\mathbb Q},\\
j_{\mathrm{gr}}&:\operatorname{Gr}^{k-1}K'(Y')_{\mathbb Q}\longrightarrow\operatorname{Gr}^kK'(X')_{\mathbb Q}.
\end{aligned}
\]
\SourcePage{471}
最后，注意 $\lambda_{-1}(F)=(-1)^{d-1}\gamma^{d-1}(F)$，所以\Term{morphism}{态射} $u^{\mathrm{gr}}$、$v^{\mathrm{gr}}$ 是从 3.7 的\Term{morphism}{态射} $u,v$ 取商得到的。故 $v^{\mathrm{gr}}\circ u^{\mathrm{gr}}=0$。

由第 VI 讲 5.8，对每个 $k$ 都有
\[
g_*(\operatorname{Fil}^{k-1}K'(Y'))\subset\operatorname{Fil}^{k-d}K'(Y),
\]
因为 $Y'=\mathbb P(\mathcal N)$，而 $\mathcal N$ 的\Term{rank}{秩}为 $d$。于是得到\Term{homomorphism}{同态}
\[
g_{\mathrm{gr}}:\operatorname{Gr}^{k-1}K'(Y')\longrightarrow\operatorname{Gr}^{k-d}K'(Y),
\]
以及与 $\mathbb Q$ 作\Term{tensor-product}{张量积}后所得的相应\Term{homomorphism}{同态}。所以 $u'^{\mathrm{gr}}$ 是由 $u'$（3.7）\Term{descent}{下降}到\Term{associated-graded}{伴随分次}所得。由于 $u'$ 是 $u$ 的\Term{left-inverse}{左逆}，$u'^{\mathrm{gr}}$ 就是 $u^{\mathrm{gr}}$ 的\Term{left-inverse}{左逆}，这同时证明了 $u^{\mathrm{gr}}$ 的\Term{injective}{单射性}。

要证明 $\operatorname{Ker}(v^{\mathrm{gr}})=\operatorname{Im}(u^{\mathrm{gr}})$，须证明：若 $y'\in\operatorname{Fil}^{k-1}K'(Y')_{\mathbb Q}$、$x\in\operatorname{Fil}^kK'(X)_{\mathbb Q}$ 满足
\[
j_*(y')+f^*(x)\in\operatorname{Fil}^{k+1}K'(X'),\tag{4.6.1}
\]
则存在 $y\in\operatorname{Fil}^{k-d}K'(Y)_{\mathbb Q}$，使
\[
\begin{aligned}
\text{i)}\quad&x+i_*(y)\in\operatorname{Fil}^{k+1}K'(X)_{\mathbb Q},\\
\text{ii)}\quad&y'-g^*(y)\lambda_{-1}(F)\in\operatorname{Fil}^kK'(Y')_{\mathbb Q}.
\end{aligned}
\]
置 $y=g_*(y')$。由第 VI 讲 5.8，$y\in\operatorname{Fil}^kK'(Y)$。再注意 $f_*$ 保持\Term{filtration}{滤过}（这里使用第 VIII 讲的结果是无妨的，因为证明该事实时并不使用 4.8）。事实上，由 1.8 ii)，$f$ 是相对\Term{virtual-dimension}{虚维数}为零的\Term{complete-intersection-morphism}{完全交态射}（第 VIII 讲 4.2）。由于 $f$ \Term{projective-morphism}{射影}，结论遂由第 VIII 讲 3.2 得到。对关系 (4.8.1) 施以 $f_*$，
\SourcePage{472}
得到
\[
f_*(j_*(y'))+f_*(f^*(x))\in\operatorname{Fil}^{k+1}K'(X)_{\mathbb Q},
\]
即由 3.6，
\[
i_*(y)+x\in\operatorname{Fil}^{k+1}K'(X)_{\mathbb Q}.
\]
为证明关系 ii)，须证
\[
y'-g^*(g_*(y'))\lambda_{-1}(F)\in\operatorname{Fil}^kK'(Y')_{\mathbb Q}.
\]
置 $z=y'-g^*(g_*(y'))\lambda_{-1}(F)$。由 3.4，
\[
\begin{aligned}
j_*(z)&=j_*(y')-j_*(g^*(g_*(y'))\lambda_{-1}(F))\\
&=j_*(y')-f^*(i_*(y)).
\end{aligned}
\]
而关系 i) 给出
\[
f^*(i_*(y))+f^*(x)\in\operatorname{Fil}^{k+1}K'(X')_{\mathbb Q}.
\]
结合 (4.6.1)，得到 $j_*(z)\in\operatorname{Fil}^{k+1}K'(X')_{\mathbb Q}$，从而
\[
j^*(j_*(z))=z(1-L)\in\operatorname{Fil}^{k+1}K'(Y')_{\mathbb Q}.
\]
以 $\bar z$ 表示 $z$ 在 $\operatorname{Gr}^{k-1}K'(Y')_{\mathbb Q}$ 中的像，以 $\overline{(1-L)}$ 表示 $1-L$ 在 $\operatorname{Gr}^1K'(Y')_{\mathbb Q}$ 中的像。于是 $\bar z\,\overline{(1-L)}=0$。此时可以应用第 VI 讲 5.11；它在与 $\mathbb Q$ 作\Term{tensor-product}{张量积}后仍然成立，因为其证明只利用 $K'(Y')$ 是\Term{free-algebra}{自由} $K'(Y)$-\Term{algebra}{代数}这一结构。故可找到 $z'\in K'(Y)_{\mathbb Q}$，使
\[
z-g^*(z')\lambda_{-1}(F)\in\operatorname{Fil}^kK'(Y')_{\mathbb Q}.
\]
施以 $g_*$，并利用 $g_*(z)=0$，得到
\[
z'\in\operatorname{Fil}^{k+1-d}K'(Y)_{\mathbb Q}.
\]
又因 $\lambda_{-1}(F)\in\operatorname{Fil}^{d-1}K'(Y')_{\mathbb Q}$，遂有
\[
z\in\operatorname{Fil}^kK'(Y')_{\mathbb Q},
\]
证明完毕。
\SourcePage{473}
现在假定 $X$ \Term{noetherian}{诺特}且 $Y$ \Term{regular-scheme}{正则}，从而 $v$ \Term{surjective}{满射}（3.7）。在 3.7 的证明中已见，对每个 $x'\in K'(X')$，存在 $y'\in K'(Y')$ 使
\[
x'=f^*(f_*(x'))+j_*(y').
\]
显然，与 $\mathbb Q$ 作\Term{tensor-product}{张量积}后仍有相同性质。因此，若 $x'\in\operatorname{Fil}^kK'(X')_{\mathbb Q}$，则存在 $(y',x)\in K'(Y')_{\mathbb Q}\times K'(X)_{\mathbb Q}$，使
\[
x'=v(y',x),\qquad x\in\operatorname{Fil}^kK'(X)_{\mathbb Q}.
\]
事实上，只须取 $x=f_*(x')$。

对 4.7 定义的\Term{ring}{环} $K'(Y')_{\mathbb Q}\times K'(X)_{\mathbb Q}$ 规定\Term{filtration}{滤过}
\[
\operatorname{Fil}^k(K'(Y')_{\mathbb Q}\times K'(X)_{\mathbb Q})
=\operatorname{Fil}^{k-1}K'(Y')_{\mathbb Q}\times\operatorname{Fil}^kK'(X)_{\mathbb Q}.
\]
由 (4.7.1) 立即可见，这是一个\Term{ring-filtration}{环滤过}。为证明对每个 $k$，\Term{homomorphism}{同态}
\[
v:\operatorname{Fil}^k(K'(Y')_{\mathbb Q}\times K'(X)_{\mathbb Q})
\longrightarrow\operatorname{Fil}^kK'(X')_{\mathbb Q}
\]
是\Term{surjective}{满射}，只须证明：若 $x'=\gamma^k(x'_1)$，且 $\varepsilon(x'_1)=0$，则存在 $(y',x)\in\operatorname{Fil}^k(K'(Y')_{\mathbb Q}\times K'(X)_{\mathbb Q})$ 使 $x'=v(y',x)$。

由前面的结果，存在 $(y'_1,x_1)\in K'(Y')_{\mathbb Q}\times K'(X)_{\mathbb Q}$ 使 $x'_1=v(y'_1,x_1)$，且 $x_1\in\operatorname{Fil}^1K'(X)_{\mathbb Q}$。由 4.7，
\[
x'=v(\gamma_*^k(y'_1,x_1)).
\]
而
\[
\gamma_*^k(y'_1,x_1)
=\left(\sum_{p=1}^k\gamma^{k-p}(x_1)\gamma^p(L,y'_1),\ \gamma^k(x_1)\right)
\]
（这是 (4.7.2) 的直接推论）。由于 $x_1$ 的\Term{augmentation}{增广}为零，$\gamma^k(x_1)\in\operatorname{Fil}^kK'(X)_{\mathbb Q}$；再由第 V 讲 5.7，
\[
\sum_{p=1}^k\gamma^{k-p}(x_1)\gamma^p(L,y'_1)\in\operatorname{Fil}^{k-1}K'(Y')_{\mathbb Q},
\]
遂得结论。
\end{proof}

\SourcePage{474}
\textbf{注记 4.9。} 在 $X$ \Term{noetherian}{诺特}、$Y$ \Term{regular-scheme}{正则}的情形下，可仿照 4.7，利用\Term{graded-ring}{分次环} $\operatorname{Gr}^*K'(X)_{\mathbb Q}$、$\operatorname{Gr}^*K'(Y)_{\mathbb Q}$，由 $i^*,i_*$ 得到的\Term{additive-homomorphism}{加法同态} $i^{\mathrm{gr}},i_{\mathrm{gr}}$，以及 $\operatorname{Gr}^*K'(Y)$ 的元素 $c^*(N)$，描述 $\operatorname{Gr}^*K'(X')_{\mathbb Q}$ 的\Term{graded-ring}{分次环}结构。同样，也可用这些数据描述由 $f^*,f_*,j^*,j_*$ \Term{descent}{下降}到\Term{associated-graded}{伴随分次}所得的\Term{homomorphism}{同态} $f^{\mathrm{gr}},f_{\mathrm{gr}},j^{\mathrm{gr}},j_{\mathrm{gr}}$。

\textbf{4.10.} 设 $X$ 为有\Term{ample-sheaf}{丰沛层}的\Term{regular-scheme}{正则概形}。已知（第 IV 讲\SGAsourcegap{原文 4.10 第一个\Term{lie-bracket}{括号}中未填第 IV 讲的具体引用编号}）用\Term{finite-locally-free}{有限秩局部自由} $\mathcal O_X$-\Term{module}{模}与用\Term{coherent-sheaf}{凝聚} $\mathcal O_X$-\Term{module}{模}构成的\Term{grothendieck-group}{Grothendieck 群}\Term{isomorphism}{同构}。利用此事实，在 $K'(X)$ 上定义一个“\Term{topological-filtration}{拓扑滤过}”：令 $\operatorname{Fil}_{\mathrm{top}}^iK'(X)$ 为\Term{support}{支集}的\Term{codimension}{余维数} $\geq i$ 的\Term{coherent-sheaf}{凝聚层}的类生成的\Term{subgroup}{子群}。容易证明，$\operatorname{Fil}_{\mathrm{top}}^iK'(X)$ 由各\Term{sheaf}{层} $\mathcal O_Z$ 的类生成，其中 $Z$ 遍历 $X$ 的\Term{codimension}{余维数} $\geq i$ 的\Term{integral-scheme}{整}\Term{closed-subscheme}{闭子概形}（第 X 讲\SGAsourcegap{原文 4.10 末尾\Term{lie-bracket}{括号}中仅有 X，具体引用编号未填}）。

\begin{SGAresult}{命题}{4.11}
设 $X$ 是有\Term{ample-sheaf}{丰沛层}的\Term{regular-scheme}{正则概形}。则对每个 $k\geq0$，有
\[
\operatorname{Fil}_{\mathrm{top}}^kK'(X)_{\mathbb Q}\subset\operatorname{Fil}_{\lambda}^kK'(X)_{\mathbb Q}.
\]
\end{SGAresult}
\textbf{附记。} 反向包含关系在不与 $\mathbb Q$ 作\Term{tensor-product}{张量积}时也成立，将在第 X 讲\SGAsourcegap{原文关于反向包含关系的引注未填第 X 讲的具体编号}中证明。
\begin{proof}
由 4.10，只须证明，若 $Z$ 为 $X$ 的\Term{codimension}{余维数}为 $k$ 的\Term{integral-scheme}{整}\Term{closed-subscheme}{闭子概形}，则
\[
\operatorname{cl}(\mathcal O_Z)\in\operatorname{Fil}_{\lambda}^kK'(X)_{\mathbb Q}.
\]
为此，对 $X$ 的\Term{integral-scheme}{整}\Term{closed-subscheme}{闭子概形}作\Term{noetherian-induction}{诺特归纳}。假定此性质对 $Z$ 的\Term{integral-scheme}{整}\Term{closed-subscheme}{闭子概形}已成立，来证明它对 $Z$ 成立。设 $\mathcal J$ 为定义 $Z$ 的 $\mathcal O_X$ 的\Term{ideal}{理想}，
\SourcePage{475}
$z$ 为 $Z$ 的\Term{generic-point}{泛点}。由于 $Z$ \Term{integral-scheme}{整}，$\mathcal J_z$ 是 $\mathcal O_{X,z}$ 的\Term{maximal-ideal}{极大理想}；又因 $X$ \Term{regular-scheme}{正则}，$\mathcal J_z$ 可由一个 $\mathcal O_{X,z}$-\Term{regular-sequence}{正则序列}生成。因此可在 $z$ 的一个\Term{neighborhood}{邻域}中找到 $\mathcal J$ 的一族\Term{section}{截面}，构成一个 $\mathcal O_X$-\Term{regular-sequence}{正则序列}。于是存在 $z$ 的一个\Term{open-neighbourhood}{开邻域} $U$，使 $Z'=X-U$ 包含于 $Z$，且 $\mathcal J|_U$ 是 $\mathcal O_U$ 的\Term{regular-ideal}{正则理想}。此时\Term{immersion}{浸入} $i:Z\cap U\to U$ 是\Term{regular-immersion}{正则浸入}，其\Term{codimension}{余维数}为 $\dim(\mathcal O_{X,z})=\operatorname{codim}(Z,X)$（EGA IV 5.1.2）。另一方面，有\Term{exact-sequence}{正合列}（第 IV 讲\SGAsourcegap{原文此处\Term{exact-sequence}{正合列}的引注未填第 IV 讲的具体编号}）
\[
K'(Z')\xrightarrow{j_*}K'(X)\xrightarrow{u^*}K'(U)\longrightarrow0,
\]
其中 $j$ 是\Term{immersion}{浸入} $Z'\to X$，$u$ 是\Term{immersion}{浸入} $U\to X$；与 $\mathbb Q$ 作\Term{tensor-product}{张量积}后得到相应的列。而 $u^*(\mathcal O_Z)=i_*(\mathcal O_{Z\cap U})$；由 4.6，$\mathcal O_Z$ 的类经 $u^*$ 所得的像 $u^*(1_Z)$ 属于 $\operatorname{Fil}_{\lambda}^kK'(U)_{\mathbb Q}$。由于 $u^*$ 是\Term{surjective}{满射} $\lambda$-\Term{lambda-homomorphism}{同态}，它诱导\Term{surjection}{满射}
\[
\operatorname{Fil}_{\lambda}^kK'(X)\longrightarrow\operatorname{Fil}_{\lambda}^kK'(U),
\]
所以 $u^*(1_Z)=u^*(x)$，其中 $x\in\operatorname{Fil}_{\lambda}^kK'(X)$；故存在 $y\in j_*(K'(Z'))$，使
\[
\operatorname{cl}(\mathcal O_Z)-y\in\operatorname{Fil}_{\lambda}^kK'(X)_{\mathbb Q}.
\]
而 $Z'\subset Z$，且因 $z\notin Z'$，$\operatorname{codim}(Z',X)>k$。于是 $x$ \Term{representable}{可表}为各\Term{sheaf}{层} $\mathcal O_{Z''_\alpha}$ 的类的组合，其中各 $Z''_\alpha$ 都是 $Z$ 的异于 $Z$ 的\Term{integral-scheme}{整}\Term{closed-subscheme}{闭子概形}。由归纳假设，$x\in\operatorname{Fil}_{\lambda}^{k+1}K'(X)_{\mathbb Q}$，证明完毕。
\end{proof}

\begin{center}\textbf{参考文献}\end{center}
\noindent [1] Cartan--Eilenberg，\textit{Homological Algebra}（《\Term{homological-algebra}{同调代数}》）。
